% Fuentes y límites de redacción: MAPA_FUENTES.md, capítulo 3.
\chapter{Datos y construcción del corpus}
\label{ch:datos}

\section{Acceso programático al BOE y adquisición}
\label{sec:adquisicion-api}
La API introducida en la sección~\ref{sec:boe-api-concepto} permite obtener
la relación de publicaciones sin recorrer manualmente la web. El proceso
parte de una fecha inicial y otra final e incluye ambos extremos. Consulta
los sumarios de forma secuencial, día a día. Se trata de un
\emph{flujo de procesamiento} o \emph{pipeline}: una secuencia en la que cada
etapa prepara la entrada de la siguiente.

El orden de adquisición importa. Primero se recuperan los metadatos del
sumario; el texto completo se descarga únicamente para los candidatos
seleccionados. El recorrido hasta la extracción puede resumirse así:

\begin{enumerate}
\item Consultar la API y obtener los sumarios del intervalo.
\item Reunir sus publicaciones y seleccionar candidatos por el título.
\item Descargar los XML de los candidatos y preparar el corpus fuente.
\item Aplicar reglas de alcance que descartan ciertos casos sin usar IA.
\item Para los casos restantes, solicitar la extracción estructurada a
Gemini, el modelo de lenguaje utilizado, y validar su respuesta.
\end{enumerate}

La selección inicial es una operación sobre títulos; no confirma todavía la
relevancia ni extrae plantas o actuaciones. Tampoco se presupone que la etapa
anterior a Gemini pueda decidir todos los casos. Las reglas y esta diferencia
se desarrollan en la sección~\ref{sec:candidatos}.

\subsection{Consulta de sumarios y recuperación del XML}
Una solicitud se dirige a una dirección del servicio, también llamada
\emph{endpoint}. El programa se comunica mediante HTTP, el protocolo que
permite solicitar recursos web. Utiliza solicitudes GET, que piden recuperar
información, a la dirección diaria:

\begin{quote}
\small
\url{https://www.boe.es/datosabiertos/api/boe/sumario/AAAAMMDD}
\end{quote}

La parte \texttt{AAAAMMDD} se sustituye por año, mes y día sin separadores.
La solicitud pide una respuesta en JSON mediante
\texttt{Accept: application/json}. JSON es un formato de datos que organiza
valores mediante nombres de campo y listas; facilita que el programa lea
cada publicación sin interpretar la apariencia de una página web.

El programa recorre las publicaciones de las secciones, departamentos y
epígrafes del sumario. Recupera el identificador BOE, la fecha, el título,
los enlaces HTML y XML y ese contexto de clasificación. HTML permite abrir
la publicación en el navegador; XML proporciona el documento que procesa
el sistema. Se comprueba que no haya identificadores duplicados.

Para descargar el XML se utiliza directamente el enlace del campo
\texttt{url\_xml} del sumario; no se construye una segunda consulta de búsqueda
por nombre de proyecto. La descarga es otra solicitud GET, separada de la
consulta JSON del sumario. El programa comprueba que la respuesta no esté
vacía y que pueda leerse como XML. Después extrae identificador, título y
fecha de los metadatos, y recorre el texto del XML en su orden documental,
con separaciones entre bloques y celdas de tablas. No se obtiene el texto
mediante reconocimiento de imágenes de un PDF.

Se comprueba que el identificador y la fecha del XML coincidan con los del
candidato. El título de la entrada se conserva del sumario; aunque se lee
también el título del XML, esta etapa no exige que ambos títulos sean iguales.
El texto preparado se reúne con sus metadatos y su procedencia.
Este resultado constituye la entrada documental de la extracción, no una
base de proyectos ya identificados.

\subsection{Conservación de la fuente}
La capa próxima a la fuente se denomina conceptualmente \emph{Bronze}.
Puede entenderse como la materia prima conservada antes de extraer la
información administrativa. En este proyecto comprende sumarios, XML y
texto preparado, junto con los datos que permiten reconocer su origen; no
significa ausencia de toda transformación, pues el XML ya se convierte a texto.

Los resultados se guardan por ejecución en una instantánea o \emph{snapshot},
es decir, un conjunto identificado de archivos correspondiente a una fase.
Incluye las respuestas de sumarios, los XML, sus metadatos de obtención y
las tablas de publicaciones, candidatos, descargas y documentos preparados.
Estas tablas utilizan Parquet, un formato de archivo para datos tabulares.

Un archivo descriptivo, denominado \emph{manifest}, declara el intervalo,
las reglas de selección, los recuentos y las identidades de los archivos.
Su función es permitir comprobar qué fuente concreta se utilizó. Las huellas
que identifican el contenido se explican en la sección~\ref{sec:identidad-documental}.
La instantánea se prepara en un directorio temporal y solo se publica cuando
supera los controles. El destino debe ser nuevo: no se sobrescribe una fuente
anterior para sustituirla por un resultado parcial.

Esta persistencia permite utilizar después documentos locales ya preparados.
No equivale a una caché automática de descargas, es decir, un mecanismo que
evite volver a pedir lo ya descargado durante otra adquisición. La etapa de
obtención actual no ofrece esa caché ni una reanudación automática desde
una descarga interrumpida. Tampoco garantiza conservar una instantánea parcial
si falla antes de publicar la fuente.

\subsection{Control de errores y límites}
El sistema distingue una consulta correcta, una respuesta de ausencia de
publicación y un fallo. En los sumarios, una respuesta 404 del servidor o del
estado BOE se registra como ausencia de publicación para esa fecha. En cambio,
un 404 al descargar un XML candidato se trata como fallo: se esperaba un
documento concreto y no se pudo recuperar.

Las solicitudes tienen un tiempo de espera de 30 segundos por defecto.
Ante ciertos fallos transitorios de conexión, tiempo de espera o respuesta
HTTP, se realizan como máximo tres reintentos adicionales, con esperas de
uno, dos y cuatro segundos. Los códigos HTTP indican el resultado de la
solicitud. Una respuesta
ilegible o una incoherencia de identidad no se resuelven aceptando el documento
como válido. Si persisten fallos de adquisición o preparación, no se publica
el corpus fuente como completo.

Estos controles ayudan a conservar una entrada comprobable, pero la API sigue
proporcionando publicaciones, no una historia de proyectos ya consolidada.
La variación de nombres y la dispersión entre documentos requieren la
agrupación posterior descrita en el capítulo~\ref{ch:metodologia}.

\pendiente{FIGURA F02 — API del BOE y construcción del corpus
\par Etiqueta prevista: \path{fig:api-corpus}. Mostrar servicio BOE, sumarios y metadatos,
filtro de títulos, XML de candidatos y documentos fuente preservados;
después, reglas de alcance con bifurcación entre negativos deterministas
y candidatos para Gemini. Mensaje: recuperar, filtrar candidatos, decidir
alcance y extraer son operaciones distintas. La descarga de XML sigue al
filtro; la relevancia no queda resuelta para todos los casos antes del modelo.
Contrastar con \texttt{boe\_source.py}, \texttt{boe\_documents.py},
\texttt{boe\_http.py}, \texttt{boe\_candidates.py} y el pipeline. Referenciar
desde esta sección después de recordar la definición de API del capítulo 2.
Usar el endpoint documentado y las URL XML del sumario, sin inventar servicios.}

\section{Identidad documental}
\label{sec:identidad-documental}
El identificador BOE permite reconocer una publicación y evitar tratarla como
varios documentos. La identidad de su contenido se controla además mediante
una \emph{huella digital} o \emph{hash}: un valor calculado a partir del
contenido que permite detectar cambios. El proyecto utiliza SHA-256, un
algoritmo concreto para calcular esas huellas.

Ambos controles responden a preguntas diferentes. El identificador indica
qué publicación es; la huella permite comprobar si la versión utilizada
coincide con la conservada. La fecha, los enlaces y las identidades de
configuración completan la procedencia. Reutilizar una extracción requiere
compatibilidad documental y de configuración; no basta con encontrar el
mismo identificador BOE.

\section{Selección de candidatos y decisión de relevancia}
\label{sec:candidatos}
El universo inicial es la relación de publicaciones recuperada mediante la
API para el intervalo consultado. La primera selección aplica una lista versionada de
expresiones sobre el \emph{título normalizado}. Normalizar significa aplicar
reglas de escritura comunes para comparar textos; aquí se homogeneizan
mayúsculas, acentos, signos y espacios. El texto fuente se conserva por separado.

La lista contiene términos de generación, como «fotovoltaica» o «parque
eólico», y expresiones administrativas más amplias, como «información
pública», «impacto ambiental» o «utilidad pública». También incluye referencias
a almacenamiento y evacuación, porque pueden formar parte de un proyecto
de generación. Este primer filtro consulta el título; ni el departamento
ni el epígrafe seleccionan por sí solos un candidato.

Se habla de \emph{candidato} porque una coincidencia textual solo indica que
el documento merece un examen posterior. Por ejemplo, la expresión
«información pública» puede aparecer en un procedimiento de carreteras.
Sería un falso positivo del filtro inicial: un documento seleccionado que no
corresponde al objeto del seguimiento. Este ejemplo describe una posibilidad
del criterio, no una observación cuantificada del corpus.

Una segunda etapa aplica reglas deterministas de alcance, es decir,
condiciones fijas que permiten resolver ciertos casos sin consultar al modelo.
Entre ellas figuran anuncios de contratación, infraestructura gasista,
almacenamiento autónomo sin planta asociada y objetos sectoriales ajenos a
la generación. Se incorporaron también reglas sobre organismos de aguas,
carreteras y costas, con las excepciones explícitas de generación e
hidroelectricidad previstas por la política vigente.

Estas reglas reducen el número de documentos que llegan a Gemini, pero no
clasifican todos los casos. Para los restantes, el modelo debe decidir el
alcance y, cuando corresponde, extraer plantas, actuaciones y otras entidades
en una respuesta estructurada. Esa respuesta se somete después a controles
deterministas y validación documental. Por tanto, no hay una confirmación
completa de relevancia de todos los candidatos antes de usar el modelo. Incluso
un título que identifica una planta y un trámite necesita comprobar la
información del documento. La relevancia exige una planta de generación
identificable y una actuación publicada vinculada a ella o a sus componentes.

Esta separación reduce trabajo innecesario, pero no demuestra que el filtro
recupere todas las publicaciones relevantes del BOE. La selección basada en
títulos y la búsqueda histórica delimitan el universo observado. La evaluación
independiente posterior tampoco mide la cobertura de los documentos descartados
por todas las etapas anteriores.

\section{Corpus de desarrollo}
El desarrollo utilizó un piloto inicial de 100 documentos y una muestra
dirigida de 40 documentos para detectar casos difíciles. Estos conjuntos
sirvieron para revisar reglas y corregir defectos; sus resultados no
constituyen una evaluación independiente del sistema final.

La referencia de desarrollo y regresión quedó fijada el 13 de agosto de
2026 sobre 140 documentos. Fijar o \emph{congelar} una versión significa
conservar una referencia identificada que no se modifica durante las fases
que dependen de ella. No implica que esa muestra sea representativa de todos
los documentos del BOE.

Además de esos conjuntos, se registraron los documentos examinados durante
auditorías y revisiones. La frontera de exposición utilizada para seleccionar
la evaluación final reúne 479 BOE. Este registro de exclusión tiene una
finalidad distinta del corpus de 140 documentos: impedir que un documento
ya utilizado para tomar decisiones de desarrollo se considere desconocido
en la evaluación. No se deben sumar ambos tamaños.

\section{Corpus final}
La estrategia final se centró en una ventana ancla de catorce días, denominada
W14, del 7 al 20 de agosto de 2026. Una ventana ancla selecciona actividad
reciente desde la que buscar publicaciones anteriores relacionadas. Tras
aplicar las reglas de alcance y excluir el documento reservado para
evaluación, el conjunto ejecutable de esta ventana contiene 48 BOE.

La recuperación retrospectiva buscó candidatos entre el 1 de enero de 2022
y el 6 de agosto de 2026. Utilizó dos niveles conservadores: coincidencia
exacta de una denominación normalizada observada, o coincidencia de términos
distintivos del nombre acompañada de tecnología o territorio compatible.
Antes de 2024 se buscó en la fuente canónica disponible; desde 2024, en los
conjuntos documentales delimitados para extracción por el plan previo P2.
La fuente canónica contiene documentos preparados tras la selección inicial,
no todas las publicaciones sin filtrar de los sumarios.

La búsqueda produjo 56 BOE históricos distintos, posteriormente extraídos y
validados. Sus enlaces eran candidatos de relación, no confirmaciones
automáticas de pertenencia al mismo proyecto. La unión de los 48 documentos
ancla y los 56 históricos constituye el corpus final de 104 publicaciones.
La versión corregida del producto contiene 86 proyectos y 80 BOE relevantes.
Por tanto, ni documento analizado, ni publicación relevante, ni proyecto
son unidades intercambiables.

El plan P2 previo planteaba una ejecución más amplia entre 2024 y agosto de
2026. Se conserva como parte del desarrollo, pero fue sustituido para el
producto final por W14 con recuperación histórica conservadora. No debe
presentarse como un corpus exhaustivo completado. Los periodos y tamaños
que resumen esta separación se recogen en la tabla~\ref{tab:corpus}.

\begin{table}[htbp]
\centering
\small
\renewcommand{\arraystretch}{1.2}
\begin{tabular}{L{3.1cm} C{1.5cm} L{4.1cm} L{4.0cm}}
\hline
Conjunto & BOE & Periodo o referencia & Función \\
\hline
Desarrollo congelado & 140 & Referencia fijada el 13/08/2026 & Desarrollo y regresión \\
Registro de exposición & 479 & Usos previos a la selección & Exclusión de evaluación \\
Ancla W14 & 48 & 07/08/2026--20/08/2026 & Actividad reciente \\
Histórico recuperado & 56 & Búsqueda: 01/01/2022--06/08/2026 & Antecedentes candidatos \\
Corpus final & 104 & Ancla más histórico & Construcción del producto \\
Evaluación independiente & 48 & Marco: 01/01/2024--20/08/2026 & Evaluación de extracción \\
\hline
\end{tabular}
\caption{Conjuntos documentales y funciones. Los intervalos de búsqueda no
son fechas mínimas y máximas observadas. Las filas no son sumables, salvo
ancla e histórico, que son disjuntos.}
\label{tab:corpus}
\end{table}

\section{Separación del conjunto de evaluación}
Un \emph{holdout} es un conjunto reservado para evaluar el sistema con
documentos que no se han utilizado para ajustar sus reglas. Permite observar
su comportamiento más allá de los ejemplos de desarrollo. En este TFM se
seleccionaron 48 BOE del marco P2 que requerían extracción mediante modelo,
después de excluir los documentos registrados como expuestos.

La selección distribuyó los documentos en seis \emph{estratos}, o grupos
definidos antes de evaluar: las series BOE-A y BOE-B combinadas con los años
2024, 2025 y 2026. Se eligieron ocho publicaciones por estrato, con semilla
fija 20260821, el valor que permite repetir la selección aleatoria. Este
diseño equilibra esos grupos; no reproduce necesariamente su proporción en
el universo del BOE.

El conjunto reservado se mantuvo separado tanto del registro de desarrollo
como de los 104 documentos del corpus W14. Su pertenencia se fijó antes de
ejecutar el sistema sobre él. La selección heredada de P2 y la construcción
posterior de W14 responden a finalidades diferentes; el holdout no es una
muestra aleatoria de los proyectos mostrados en la aplicación.

La anotación humana, la fijación de las reglas de evaluación y las medidas
utilizadas se explicarán en el capítulo~\ref{ch:evaluacion}. La tabla anterior
describe los conjuntos y no anticipa resultados de rendimiento.
